\documentclass[11pt]{article}
\usepackage[review]{acl}
\usepackage{times}
\usepackage{latexsym}
\usepackage[T1]{fontenc}
\usepackage[utf8]{inputenc}
\usepackage{microtype}
\usepackage{inconsolata}
\usepackage{graphicx}
\usepackage{amsmath,amssymb,amsthm}
\usepackage{booktabs}
\usepackage{multirow}
\usepackage{array}
\usepackage{xcolor}
\usepackage{algorithm}
\usepackage{algpseudocode}
\usepackage{url}

\newtheorem{proposition}{Proposition}

\title{CAP-Diag: Diagnosing Detector-Statistic-Accessible Information
in LLM Watermarks under Controlled Rewrite}

\author{Anonymous Submission \\
  ACL 2027 (anonymized) \\
  \texttt{anonymous@example.org}}

\begin{document}
\maketitle

\begin{abstract}
Watermark detectors are evaluated by AUROC / TPR / robustness benchmarks, but
those metrics do not say how much information the chosen detection statistic
$S$ still carries about the watermark flag $W$ once the text has been
meaning-preservingly rewritten. We propose \textbf{CAP-Diag}, a diagnostic
that estimates the \emph{detector-statistic-accessible} MI $\hat I(W;S)$
under a controlled rewrite ladder (R0/R1/R1.5/R2), with a Fano-style
chosen-statistic error floor for any classifier consuming $S$
(Prop.~\ref{prop:bayes}). On $200$ synthetic and $300$ Wikipedia sources
across three token-level watermark families (KGW, UW, EXP/Gumbel),
CAP-Diag exposes a \emph{regime/length-coupled low-information bottleneck}:
source-paired R1/R0 MI ratio $\approx 4$--$15\times$ across families and
corpora, KGW/UW Wilcoxon $p<10^{-34}$. A controlled length-matched re-analysis (R1-short generation at target
$30$--$50$ tok, mean $\approx 34$ tok matching R0) shows that, at matched
length, source-paired R1-short / R0 ratio is $1.57\times$ [$1.21, 2.30$]
for KGW ($n_{\mathrm{pairs}}=148$) and $2.71\times$ [$1.74, 4.52$] for UW
($n_{\mathrm{pairs}}=111$); both CIs exclude $1$. Length is the dominant
driver but a regime-specific residual component is supported with
adequate power across both families. The bottleneck drives a \emph{Low-Capacity Rewrite Attack}
that re-rewrites watermarked R2 text into the diagnosed R0 regime
($\approx 80\%$ MI drop), and \textbf{CAP-Triage} converts the diagnosis
into an abstention readout with held-out source and held-out
Wikipedia-domain validation. Mechanism boundaries are bounded references:
a SemStamp-style k-LSH reference also exhibits R0$<$R1, while a
POSTMARK-style content-insertion sanity check does not (R1/R0
$\approx 1.0$). We complement the long-text reliability lens
of~\citet{kirchenbauer2024reliability}: long spans are reliably detectable;
CAP-Diag characterises the short, compressed, fact-skeleton regime in
which $S$ no longer carries enough information. CAP-Diag is
statistic-conditional, not Shannon capacity; not a benchmark, not a new
watermark, not a defense.
\end{abstract}

\section{Introduction}
\label{sec:intro}

LLM watermarks~\citep{kirchenbauer2023watermark,zhao2024provable,kuditipudi2024robust}
underpin provenance pipelines that must remain readable after the watermarked
text is rewritten, summarised, or otherwise transformed by downstream users.
Standard evaluation reports detection metrics (AUROC, TPR at fixed FPR) and
attack-robustness benchmarks~\citep{piet2024markmywords,tu2024waterbench,liang2024waterpark};
recent reliability work~\citep{kirchenbauer2024reliability} shows that
sufficiently long observed watermarked spans remain reliably detectable
even under substantial paraphrase or human edits. \emph{What these metrics
do not answer is how much information the chosen detector statistic $S$
itself still carries about the watermark flag $W$ once the text has been
rewritten under a controlled meaning-preserving transformation.} A detector
that scores AUROC$=0.99$ on long high-strength text may still face a regime
in which $S$ has near-zero information about $W$ --- and AUROC alone does
not say which regime that is, or what to do at deployment when it is.

\paragraph{Diagnostic object.}
We define and operationalise the chosen-statistic-conditional mutual
information $\hat I(W;S)$ under a controlled rewrite ladder: regimes R0
(short fact-skeleton), R1 (short paraphrase), R1.5 (looser surface), R2
(long elaboration), each gated by NLI bidirectional non-contradiction,
fact-skeleton match, and length band. Prop.~\ref{prop:bayes} converts
small $\hat I(W;S)$ into a Fano-style lower bound on the error rate of
\emph{any} classifier that consumes $S$, including the Bayes-optimal one.
This is statistic-conditional --- not Shannon capacity, not an
all-detector upper bound.

\paragraph{Diagnostic protocol.}
\textbf{CAP-Diag} packages the measurement as a reusable
input/output contract (Sec.~\ref{sec:capdiag}, Table~\ref{tab:contract}):
inputs are a source set, regime ladder, utility gate, watermark family,
detector statistic $S$, and prior $\pi$; the procedure runs a clipped
Miller--Madow estimator on a balanced sample with source-paired bootstrap,
AUROC/TPR alongside, length and permutation nulls, and a risk-coverage
readout; the output is an accessible-MI ladder, a Fano floor, AUROC/TPR
comparison, confound diagnostics, the Low-Capacity Rewrite Attack
consequence, and a CAP-Triage abstention recommendation.

\paragraph{Main finding: regime/length-coupled low-information bottleneck.}
Across three token-level families (KGW/UW/EXP) on synthetic and
naturalistic English, fact-skeleton R0 rewrites carry markedly less
detector-accessible information than R1 paraphrases (source-paired R1/R0
MI ratio $\approx 4$--$15\times$, KGW/UW Wilcoxon $p<10^{-34}$). Length is
\emph{not} a caveat: a length-only baseline classifies R0 vs R1 at
AUROC$=0.996$, and we adopt the length-as-sample-complexity lens
of~\citet{kirchenbauer2024reliability} to position the finding as
\emph{regime/length-coupled} rather than length-orthogonal. A controlled
length-matched re-analysis (Sec.~\ref{sec:length_disent}; R1-short prompt
at target $30$--$50$ tokens, mean $\approx 34$ matching R0) shows that
when length is matched, source-paired R1-short / R0 MI ratio is
$1.57\times$ [$1.21, 2.30$] for KGW ($n_{\mathrm{pairs}}=148$) and
$2.71\times$ [$1.74, 4.52$] for UW ($n_{\mathrm{pairs}}=111$); both CIs
exclude $1$ with adequate power across both families. The bottleneck is
therefore length-coupled \emph{but not length-explained}: a residual
regime-specific component remains under controlled length matching.

\paragraph{Consequence and readout.}
Because R0 is a diagnosed low-information regime, an attacker can
re-rewrite watermarked R2 text into R0-style fact skeleton without changing
key facts; the resulting MI drop is $81\%$ (KGW) / $77\%$ (UW) with
bootstrap CIs (\textbf{Low-Capacity Rewrite Attack}, Sec.~\ref{sec:lowcap_attack}).
\textbf{CAP-Triage} (Sec.~\ref{sec:cap_triage}) converts the same
low-information diagnosis into an abstention readout that reduces
false-confident emissions on R0-dominated documents using only an
observable length proxy, with held-out source and held-out Wikipedia-domain
validation. CAP-Triage is a deployment readout, not a defense: it does not
recover lost information, it tells the deployer when not to emit.

\paragraph{Mechanism boundaries.}
A SemStamp-style k-LSH sentence-rejection reference implementation also
exhibits the R0$<$R1 pattern (Sec.~\ref{sec:semstamp_neg}); a POSTMARK-style
content-insertion sanity check (with simulated rewrite-survival, not a
faithful reproduction of POSTMARK) does \emph{not} (R1/R0 ratio
$\approx 1.0$, Sec.~\ref{sec:postmark}). We read these as bounded
\emph{reference contrasts} rather than category-level evidence: they
suggest the bottleneck is tied to token-distribution-deviation statistics
and is not universal across watermark mechanism classes, but they are not a
mechanism-class theorem.

\paragraph{Contributions.}
\begin{enumerate}\itemsep0pt
\item \textbf{Diagnostic object and protocol.} Detector-statistic-accessible
MI $\hat I(W;S)$ with a Fano-style chosen-statistic error floor
(Prop.~\ref{prop:bayes}), packaged as the CAP-Diag protocol with an explicit
input/output contract (Table~\ref{tab:contract}).
\item \textbf{Regime/length-coupled low-information bottleneck.} Source-paired
R1/R0 MI ratio $\approx 4$--$15\times$ across three token-level families on
two corpora; length is part of the finding, with a length-band-matched
re-analysis confirming a residual regime-specific component
(Sec.~\ref{sec:confound}--\ref{sec:length_disent}).
\item \textbf{Low-Capacity Rewrite Attack as diagnostic consequence.} A
meaning-preserving R0-style re-rewrite of watermarked R2 text drops MI by
$\approx 80\%$ (KGW $81.4\%$, UW $77.1\%$ with bootstrap CIs).
\item \textbf{CAP-Triage abstention readout.} Risk-coverage / abstention
diagnostic with held-out source and held-out Wikipedia-domain
validation; not a defense.
\end{enumerate}

Supporting validation (mechanism-boundary references SemStamp-style and
POSTMARK-style; LLM-judge gate validation; MI bin-count, classifier-LB,
gate-NLI, secret-key, skeleton-policy and adaptive-detector sensitivities;
permutation nulls; per-domain analysis; three honest negative findings) is
reported in Sec.~\ref{sec:results} and the appendices, but is not part of
the main contribution list above.

\begin{figure*}[t]
    \centering
    \includegraphics[width=0.97\textwidth]{../figures/fig1_capdiag.pdf}
    \caption{CAP-Diag protocol. Source passages are rewritten under each
    regime $R$ via a watermarked generator $\mathcal{W}_f(\theta)$. The
    utility-preserving gate $\mathcal{G}$ enforces NLI bidirectional
    non-contradiction, fact-skeleton match, and length band. Surviving
    drafts feed a Miller--Madow MI estimator with source-paired bootstrap.
    Length-matched and permutation nulls test whether observed regime
    differences reflect more than length or regime label. An abstention rule
    converts accessible-information estimates into deployment behaviour.}
    \label{fig:capdiag}
\end{figure*}

\section{Background and Setup}
\label{sec:background}

\subsection{Watermark Families}
\textbf{KGW}~\citep{kirchenbauer2023watermark} hashes the previous token
with the secret to derive a per-step green list of $\gamma V$ tokens and
adds logit bias $\delta$ on green tokens.
\textbf{Unigram-Watermark (UW)}~\citep{zhao2024provable} uses a single
fixed green list seeded by the secret. Both share the $z$-score detector
$z=(n_{\mathrm{green}}-\gamma T)/\sqrt{T\gamma(1-\gamma)}$.
\textbf{EXP/Gumbel}~\citep{aaronson2023watermark,kuditipudi2024robust}
replaces sampling with $x_t = \arg\max_v \log r_v / p_v$ where $r_v$ is
hash-derived from context plus secret; the detector aggregates
$S_T = \sum_t -\log(1-r_{x_t})$. We also evaluate a
\textbf{SemStamp-style k-LSH} sentence-rejection
\emph{reference implementation}~\citep{hou2024semstamp} as a non-token-level
reference (not the official SemStamp pipeline, not a category-level
representative); the detector reports the fraction of generated sentences whose SBERT
embedding falls in a secret-keyed green bucket. We use $\gamma=0.25$ for
greenlist methods throughout.

\subsection{Rewrite Regimes}
We study four utility-preserving rewrite regimes:
\textbf{R0} fact-skeleton rephrase ($\sim$32 synth tok / $\sim$60 nat tok);
\textbf{R1} short NLI-preserved paraphrase ($\sim$50 / $\sim$86 tok);
\textbf{R1.5} looser-form NLI-preserved paraphrase ($\sim$46 tok);
\textbf{R2} longer NLI-preserved rewrite ($\sim$103 / $\sim$160 tok).
Prompts in Appendix~\ref{app:prompts}. Examples in Appendix~\ref{app:examples}.

\subsection{Detector-Statistic-Accessible MI}
Treat $W\in\{0,1\}$ as ``is this draft watermarked?'' and $S$ as the
detector statistic. With pooled samples $z^{(\mathrm{wm})}\!\cup\!
z^{(\mathrm{no})}$ and $\pi_{\mathrm{wm}}$ the watermarked fraction, we
estimate
\begin{align*}
\hat{I}(W;S) =\; & \hat H_{\mathrm{MM}}(\mathrm{pooled}) \\
              & - \pi_{\mathrm{wm}} \hat H_{\mathrm{MM}}(z^{(\mathrm{wm})}) \\
              & - (1{-}\pi_{\mathrm{wm}}) \hat H_{\mathrm{MM}}(z^{(\mathrm{no})}),
\end{align*}
where $\hat H_{\mathrm{MM}}$ is the Miller--Madow-corrected histogram
entropy on $20$ equal-width bins over the pooled range. We clip each
bootstrap replicate to $[0,1]$ before percentile computation, so CIs lie in
$[0,1]$ by construction. Throughout we refer to $\hat I(W;S)$ as
\emph{detector-statistic-accessible MI}; we use ``capacity'' only as
shorthand after the formal statement below.

\begin{proposition}[Detection-error lower bound under $I(W;S)$]
\label{prop:bayes}
Let $W\in\{0,1\}$ with prior $\pi$ and let $S$ be any (possibly randomized)
function of the observed text. For \emph{any} (possibly randomized)
classifier $\hat W = h(S)$ that consumes exactly $S$, denote its error
$P_e = \Pr[\hat W \ne W]$. Fano's inequality gives
\[
H(W \mid S) \le H_b(P_e) + P_e \cdot \log_2(|\mathcal{W}|-1),
\]
which for binary $W$ reduces to $H(W) - I(W;S) \le H_b(P_e)$, hence
\[
P_e \ge H_b^{-1}\!\bigl(H(W) - I(W;S)\bigr),
\]
where $H_b^{-1}$ denotes the branch on $[0,1/2]$ (and the bound is
trivial if $P_e > 1/2$). The bound holds in particular for the
Bayes-optimal classifier $\hat W^\star$, whose error rate $P_e^\star$ is
the minimum over all $h$. When $I(W;S) \to 0$, $H(W\mid S)\to H(W)$ so
the lower bound forces $P_e \ge H_b^{-1}(H(W))$ for \emph{every} $h$, and
the Bayes-optimal classifier saturates the prior-only baseline
$P_e^\star \to \min(\pi, 1-\pi)$. Sub-optimal classifiers may of course
have $P_e$ strictly larger than the prior baseline.
\end{proposition}

Prop.~\ref{prop:bayes} formalizes what CAP-Diag measures: detection error
floor for \emph{any} classifier that consumes the chosen $S$. We do
\emph{not} claim that richer statistics $S'\supsetneq S$ obey the same
floor; testing stronger statistics is the purpose of our adaptive-detector
sensitivity (Sec.~\ref{sec:robust}).

\paragraph{Prior $\pi$ and prior-shift interpretation.}
Throughout this paper we use the empirical balanced prior $\pi{=}0.5$
(equal wm and no-wm samples per cell). The MI quantity $I(W;S)$ depends on
the joint distribution and hence on the prior; in deployment under prior
shift to $\pi'$, one would re-weight to obtain $I_{\pi'}(W;S)$ for use in
the Fano floor below, or interpret our reported $I_{0.5}(W;S)$ as the MI
under a balanced detection scenario. Under deployment prior $\pi'{=}0.1$,
$H(W)\!=\!H_b(0.1)\!\approx\!0.469$ bit, which loosens the bound
$P_e \ge H_b^{-1}(H(W)-I(W;S))$ proportionally; the qualitative direction
is unchanged: low $\hat I(W;S)$ forces the Bayes-optimal error toward the
prior-baseline $\min(\pi',1{-}\pi')$, which for $\pi'{=}0.1$ is $0.1$.

\paragraph{Classifier lower bound details.}
The classifier lower bound on $I(W;S)$ used in Table~\ref{tab:mi_robust} is
a cross-entropy lower bound on the balanced sample ($n_{\text{wm}}{=}n_{\text{no}}$,
$\pi{=}0.5$): train a logistic classifier $\hat p_\theta(W{=}1\mid S)$ via
$5$-fold stratified CV; compute
$\hat I_{\text{LB}} = 1 - \mathrm{NLL}_{\text{CV}}(\hat p_\theta)$
where NLL is the cross-entropy in bits. This is a strict lower bound on
$I(W;S)$ for the given feature vector under the balanced prior. The base
feature is the per-text $z$-score $z$; the $S'$ variant adds
$n_{\text{tokens}}$, green count, and green rate.

\paragraph{Defining $\text{ppl\_ratio}$.} The PPL ratio used in
Sec.~\ref{sec:c2} is the median per-source ratio of the watermarked
generator's perplexity on a rewrite, divided by the same generator's
perplexity on the original source, both computed by Qwen2.5-7B-Instruct.

\paragraph{Utility-Preserving Gate.}
\label{sec:gate}
Each rewrite must pass: (i) NLI bidirectional non-contradiction
(cross-encoder DeBERTa-v3) with $P(\text{contradiction})<0.5$ in both
directions; (ii) fact-skeleton match: required fact keys appear verbatim
modulo date/time variants; (iii) length within the regime's band.
Gate sensitivity (NLI threshold $\{0.3, 0.5, 0.7\}$ + skeleton policy
$\{$strict, relaxed, soft$\}$) is reported in Sec.~\ref{sec:robust}.

\paragraph{Source-paired Bootstrap.}
$B=500$ source-paired replicates (seed $20260513$) resampled by
\emph{(item\_id, agent\_seed)} are used for the canonical
Tables~\ref{tab:main_results}, \ref{tab:mi_vs_det}, \ref{tab:mi_robust}
and~\ref{tab:gate_sens}; a small number of earlier-round MI-estimator
sensitivity cells in App.~\ref{app:mi_sens} used $B=1000$.
Permutation $p$ uses the exact convention $p=(b+1)/(B+1)$.

\section{The CAP-Diag Protocol}
\label{sec:capdiag}

CAP-Diag is a reusable diagnostic protocol, not a metric and not a
benchmark. Table~\ref{tab:contract} states its input/output contract.

\begin{table}[t]
\centering
\small
\caption{CAP-Diag input/output contract.}
\label{tab:contract}
\setlength{\tabcolsep}{4pt}
\begin{tabular}{p{0.95\columnwidth}}
\toprule
\textbf{Input.} (a) Source set $\mathcal{D}=\{x_i\}$; (b) rewrite regime
ladder $R\in\{\mathrm{R0},\mathrm{R1},\mathrm{R1.5},\mathrm{R2}\}$ defined
by prompt + length band; (c) utility gate $\mathcal{G}$ (NLI bidirectional
non-contradiction + fact-skeleton match + length band); (d) watermark family
$f$ and strength $\theta$ (defining a generator $\mathcal{W}_f(\theta)$);
(e) chosen detector statistic $S$ (e.g. KGW/UW $z$, EXP/Gumbel
$\sum -\log(1-r)$); (f) empirical prior $\pi$ over $W\in\{0,1\}$. \\
\midrule
\textbf{Procedure.} (1) For each $(i, r, f, \theta)$ sample
$\mathcal{W}_f(\theta)(\pi_r(x_i))$ until $K$ drafts pass $\mathcal{G}$;
(2) compute $\hat I(W;S)$ via clipped Miller--Madow on a balanced
$n_{\mathrm{wm}}{=}n_{\mathrm{no}}$ sample at the stated $\pi$; (3) compute
AUROC / TPR\,@\,FPR on the same balanced cells; (4) source-paired bootstrap
$B$ replicates; (5) length-bin and regime-label permutation nulls; (6)
risk-coverage readout via an observable proxy (length or
source-relative PPL). \\
\midrule
\textbf{Output.} (a) Accessible-information ladder $\hat I_r(W;S)$ with
$95\%$ CIs; (b) Fano-style chosen-statistic error floor
$P_e \ge H_b^{-1}(H(W){-}\hat I(W;S))$ for any classifier consuming $S$
(Prop.~\ref{prop:bayes}); (c) AUROC/TPR comparison; (d) confound
diagnostics (length, key density, entropy, source-paired Wilcoxon); (e) a
Low-Capacity Rewrite Attack consequence quantified as MI drop with bootstrap
CI; (f) a CAP-Triage risk-coverage / abstention recommendation. \\
\midrule
\textbf{Scope.} Statistic-conditional and prior-conditional only. Not
Shannon channel capacity. Not an all-detector upper bound. Not a benchmark.
Not a new watermark method. Not a defense. CAP-Triage is an abstention
readout, not a detector.\\
\bottomrule
\end{tabular}
\end{table}

The full algorithm is in Appendix~\ref{app:algorithm} and a diagrammatic
view is in Fig.~\ref{fig:capdiag}.

Three properties distinguish CAP-Diag from prior watermark evaluation
lines (Sec.~\ref{sec:related}, Table~\ref{tab:related}):
\textbf{(i) metric}: MarkMyWords~\citep{piet2024markmywords},
WaterBench~\citep{tu2024waterbench}, and WaterPark~\citep{liang2024waterpark}
report combinations of AUROC, TPR at fixed FPR, and (for WaterPark) attack
robustness; CAP-Diag reports $\hat I(W;S)$ together with the standard
metrics.
\textbf{(ii) controls}: a pre-registered ladder of rewrite regimes,
length-matched buckets with regime-specific baselines, permutation tests
over regime labels.
\textbf{(iii) output}: an abstention rule that says \emph{when not to
emit} a detector decision.

\section{Experimental Setup}
\label{sec:setup}

\textbf{Synthetic corpus.} $200$ fact-skeleton episodes (entity, value,
date), 2 source variants A/B, 7-strength sweep $\{0,0.5,1,1.5,2,3,5\}$ for
KGW/UW; EXP/Gumbel at peak.
\textbf{Naturalistic corpus.} $300$ Wikipedia REST passages across $15$
content domains; each $\geq\!1$ named entity + $\geq\!1$ date/numeric
(spaCy NER); R0/R1 cover all $300$, R1.5/R2 cover a $150$-source subset.
\textbf{Models.} Qwen2.5-7B-Instruct~\citep{qwen2_5_7b} primary;
Llama-3.1-8B-Instruct~\citep{grattafiori2024llama3} sanity check
(App.~\ref{app:llama}).
\textbf{NLI gate.} cross-encoder/nli-deberta-v3-base.
\textbf{Reproducibility.} Three additional secret keys
$\{13371337, 99999991, 42424242\}$. Prompts, source IDs, and raw scores
will be released in the supplementary upon publication.

\section{Results}
\label{sec:results}

We organise the empirical evidence around six research questions, each
serving one named claim from the contribution list:

\begin{itemize}\itemsep0pt
\item \textbf{RQ1} (\S\ref{sec:c1}): does controlled rewrite create an
accessible-information ladder $\hat I_{\mathrm{R0}} < \hat I_{\mathrm{R1}}
< \hat I_{\mathrm{R2}}$ across watermark families and corpora?
\item \textbf{RQ2} (\S\ref{sec:c1}, \emph{MI vs detection metrics}): what
does $\hat I(W;S)$ diagnose beyond AUROC/TPR?
\item \textbf{RQ3} (\S\ref{sec:robust}--\S\ref{sec:c2}): what mechanism
explains the bottleneck, and how is it length/regime-coupled?
\item \textbf{RQ4} (\S\ref{sec:lowcap_attack}): can the diagnosed
low-information regime be actively triggered as an attack consequence?
\item \textbf{RQ5} (\S\ref{sec:cap_triage}): can the diagnostic be
turned into a deployment-side abstention readout?
\item \textbf{RQ6} (\S\ref{sec:semstamp_neg}, \S\ref{sec:postmark}): what
are the mechanism boundaries --- which watermark mechanism classes
exhibit the bottleneck and which do not?
\end{itemize}

\subsection{RQ1: Accessible-Information Ladder Across Regimes}
\label{sec:c1}

\begin{table*}[t]
\centering
\small
\caption{Accessible-information ladder (peak strength). Single-bit
detector-statistic-accessible MI in bits with $95\%$ source-paired
bootstrap CIs (clipped to $[0,1]$).
\textbf{Tables 1 and 2 share the same canonical sample, estimator, and bootstrap}
(Round-7 unification): clipped Miller--Madow with $20$ equal-width bins,
$B{=}500$ source-paired bootstrap, seed $20260513$. Empirical prior
$\pi{=}0.5$: per cell we use $n_{\text{wm}}{=}n_{\text{no}}$ (synthetic
$200/200$, naturalistic R0/R1 $300/300$, naturalistic R1.5/R2 subset
$150/150$); larger pools are downsampled with a fixed seed to maintain
$\pi{=}0.5$. ``mean tok'' is mean number of generated tokens after the
utility gate.}
\label{tab:main_results}
\begin{tabular}{l l c c c c c c}
\toprule
& & \multicolumn{3}{c}{Synthetic} & \multicolumn{3}{c}{Naturalistic} \\
\cmidrule(r){3-5}\cmidrule(l){6-8}
Regime & Family & MI & 95\% CI & mean tok & MI & 95\% CI & mean tok \\
\midrule
\multirow{3}{*}{R0} & KGW & 0.131 & [0.103, 0.232] & 33 & 0.207 & [0.179, 0.283] & 60 \\
                    & UW  & 0.136 & [0.107, 0.221] & 32 & 0.136 & [0.107, 0.202] & 61 \\
                    & EXP & 0.074 & [0.062, 0.158] & 32 & 0.058 & [0.046, 0.118] & 61 \\
\midrule
\multirow{3}{*}{R1} & KGW & 0.882 & [0.834, 0.942] & 52 & 0.858 & [0.818, 0.897] & 87 \\
                    & UW  & 0.948 & [0.920, 0.985] & 50 & 0.724 & [0.682, 0.791] & 86 \\
                    & EXP & 0.751 & [0.701, 0.842] & 50 & 0.689 & [0.637, 0.751] & 89 \\
\midrule
\multirow{3}{*}{R1.5} & KGW & 0.850 & [0.796, 0.917] & 48 & 0.861$^\dagger$ & [0.806, 0.938] & 104 \\
                      & UW  & 0.826 & [0.789, 0.897] & 48 & 0.777$^\dagger$ & [0.728, 0.883] & 95 \\
                      & EXP & 0.488 & [0.419, 0.600] & 44 & --- & --- & --- \\
\midrule
\multirow{3}{*}{R2}   & KGW & 0.995 & [0.988, 1.000] & 101 & 1.000$^\dagger$ & [1.000, 1.000] & 173 \\
                      & UW  & 0.985 & [0.965, 1.000] & 113 & 0.954$^\dagger$ & [0.929, 0.987] & 155 \\
                      & EXP & 0.977 & [0.960, 1.000] & 142 & --- & --- & --- \\
\bottomrule
\end{tabular}
\\[2pt]
{\footnotesize $^\dagger$ Naturalistic R1.5/R2 reported on the $150$-source
subset of the $300$-source corpus; the EXP family was not extended to
naturalistic R1.5/R2 in this round (token-level KGW/UW only).}
\end{table*}


Table~\ref{tab:main_results} gives peak-strength MI per (regime, family,
corpus). \textbf{Tables 1 and 2 share an identical balanced
sample/estimator/bootstrap} (clipped Miller--Madow, $20$ bins, $B{=}500$
source-paired bootstrap, peak $\delta{=}5$ or $\beta{=}1$, empirical prior
$\pi{=}0.5$). Synthetic uses $n_{\text{wm}}{=}n_{\text{no}}{=}200$
(variant A; larger no-wm pools downsampled with a fixed seed to maintain
balance); naturalistic R0/R1 uses $n_{\text{wm}}{=}n_{\text{no}}{=}300$;
nat R1.5/R2 use a $150$-source subset. R0 $<$ R1 $<$ R2 holds for KGW,
UW, and EXP/Gumbel on the synthetic corpus; for naturalistic text the
KGW/UW four-regime ladder holds on the $150$-source subset, while
EXP/Gumbel naturalistic is reported on R0/R1 only (this round does not
extend EXP to nat R1.5/R2). The SemStamp-style k-LSH reference also shows
R0$<$R1 (Sec.~\ref{sec:semstamp_neg}); the POSTMARK-style reference does
not (Sec.~\ref{sec:postmark}).

\paragraph{R1/R0 ratios.} Source-paired bootstrap-inside-loop on the
unified sample: synthetic KGW $4.14\times$ [$2.95,4.81$], UW
$4.80\times$ [$3.27,5.55$], EXP $14.89\times$ [$6.80,18.51$];
naturalistic KGW $4.14\times$ [$3.23,4.45$], UW $5.32\times$
[$3.96,5.95$], EXP $11.84\times$ [$6.29,14.71$].
All $p_{\mathrm{exact}}{=}5{\times}10^{-4}$. Ratios computed on point
estimates would inflate the gap (Jensen on division near zero).

\paragraph{Naturalistic R1.5/R2.} The full $300$-source naturalistic corpus
covers R0/R1; for the four-regime ladder we evaluate a $150$-source
subset of the same corpus on KGW/UW: KGW R0/R1/R1.5/R2$=0.21/0.86/0.86/1.00$;
UW $=0.14/0.72/0.78/0.95$ (R0/R1 from full $300$-source corpus, R1.5/R2
from the $150$-source subset; see Table~\ref{tab:main_results}).
EXP/Gumbel was not extended to nat R1.5/R2 in this round.

\paragraph{R1.5 is a more-selected, looser-paraphrase cell.}
The utility gate accepts R0 at $0.955$, R1 at $0.762$, R2 at $0.830$, but
\textbf{R1.5 at $0.297$} (Appendix~\ref{app:gate}). R1.5 results in
Table~\ref{tab:main_results} and Table~\ref{tab:mi_vs_det} should be
interpreted as a more selected pool of looser-form paraphrases rather than
as a regime fully comparable to R0/R1/R2; the looser-form R1.5 prompt
generates more NLI-rejected drafts, and the surviving R1.5 drafts are
those that did pass the bidirectional non-contradiction check. This does
not affect the headline R0$<$R1 contrast (both have high pass rates) but it
qualifies any reading of the full ladder as a continuous strength sweep.

\paragraph{RQ2 --- MI vs standard detection metrics.}
\begin{table*}[t]
\centering
\small
\caption{MI vs standard detection metrics on the same cells as
Table~\ref{tab:main_results} (\emph{identical balanced sample, estimator,
and bootstrap}; Round-7 unification, empirical prior $\pi{=}0.5$).
KGW/UW AUROC saturates above $0.97$ for R1, R1.5, and R2 on both corpora;
EXP/Gumbel AUROC saturates at R1 ($0.98$) and R2 ($1.00$) but reaches
only $0.93$ at R1.5. The R0 row is not a case of AUROC hiding MI
structure: AUROC there is also low ($\leq 0.80$) and TPR@$5\%$FPR
$\leq 0.31$. The MI complement is most informative in the
high-AUROC saturation regime: at R1/R1.5/R2, AUROC ranges $0.92$--$1.00$
while MI spans $0.49$--$1.00$, revealing residual accessible-information
structure that AUROC compresses to its ceiling. Computed from $B{=}500$
source-paired bootstrap.}
\label{tab:mi_vs_det}
\setlength{\tabcolsep}{4pt}
\begin{tabular}{l l l c c c c}
\toprule
Corpus & Regime & Family & MI (CI) & AUROC (CI) & TPR@$1\%$ & TPR@$5\%$ \\
\midrule
\multirow{12}{*}{syn} & \multirow{3}{*}{R0}   & KGW & 0.131 [0.103, 0.232] & 0.739 [0.690, 0.784] & 0.095 & 0.225 \\
                      &                        & UW  & 0.136 [0.107, 0.221] & 0.711 [0.661, 0.760] & 0.065 & 0.140 \\
                      &                        & EXP & 0.074 [0.062, 0.158] & 0.641 [0.589, 0.693] & 0.145 & 0.230 \\
\cmidrule(l){2-7}
                      & \multirow{3}{*}{R1}   & KGW & 0.882 [0.834, 0.942] & 0.997 [0.994, 0.999] & 0.935 & 0.975 \\
                      &                        & UW  & 0.948 [0.920, 0.985] & 0.999 [0.997, 1.000] & 0.980 & 0.990 \\
                      &                        & EXP & 0.751 [0.701, 0.842] & 0.982 [0.969, 0.991] & 0.705 & 0.880 \\
\cmidrule(l){2-7}
                      & \multirow{3}{*}{R1.5} & KGW & 0.850 [0.796, 0.917] & 0.991 [0.980, 0.997] & 0.930 & 0.965 \\
                      &                        & UW  & 0.826 [0.789, 0.897] & 0.992 [0.986, 0.997] & 0.870 & 0.955 \\
                      &                        & EXP & 0.488 [0.419, 0.600] & 0.925 [0.897, 0.949] & 0.350 & 0.665 \\
\cmidrule(l){2-7}
                      & \multirow{3}{*}{R2}   & KGW & 0.995 [0.988, 1.000] & 1.000 & 1.000 & 1.000 \\
                      &                        & UW  & 0.985 [0.965, 1.000] & 1.000 & 0.995 & 1.000 \\
                      &                        & EXP & 0.977 [0.960, 1.000] & 1.000 & 0.995 & 1.000 \\
\midrule
\multirow{8}{*}{nat}  & \multirow{3}{*}{R0}   & KGW & 0.207 [0.179, 0.283] & 0.796 [0.762, 0.831] & 0.163 & 0.310 \\
                      &                        & UW  & 0.136 [0.107, 0.202] & 0.736 [0.697, 0.776] & 0.117 & 0.277 \\
                      &                        & EXP & 0.058 [0.046, 0.118] & 0.656 [0.616, 0.702] & 0.060 & 0.180 \\
\cmidrule(l){2-7}
                      & \multirow{3}{*}{R1}   & KGW & 0.858 [0.818, 0.897] & 0.993 [0.989, 0.997] & 0.860 & 0.963 \\
                      &                        & UW  & 0.724 [0.682, 0.791] & 0.977 [0.968, 0.986] & 0.773 & 0.890 \\
                      &                        & EXP & 0.689 [0.637, 0.751] & 0.971 [0.960, 0.981] & 0.667 & 0.860 \\
\cmidrule(l){2-7}
                      & \multirow{2}{*}{R1.5$^\dagger$} & KGW & 0.861 [0.806, 0.938] & 0.994 [0.986, 0.999] & 0.900 & 0.973 \\
                      &                                  & UW  & 0.777 [0.728, 0.883] & 0.990 [0.980, 0.996] & 0.860 & 0.947 \\
\cmidrule(l){2-7}
                      & \multirow{2}{*}{R2$^\dagger$}   & KGW & 1.000 & 1.000 & 1.000 & 1.000 \\
                      &                                  & UW  & 0.954 [0.929, 0.987] & 0.999 [0.998, 1.000] & 0.967 & 1.000 \\
\bottomrule
\end{tabular}
\\[2pt]
{\footnotesize $^\dagger$ $150$-source subset; EXP not extended to nat R1.5/R2.}
\end{table*}

Table~\ref{tab:mi_vs_det} reports AUROC and TPR@$1\%/5\%$FPR on identical
cells. The argument is \emph{two-layered}: \textbf{at R0}, AUROC is also
low ($0.64$--$0.80$) and TPR@$5\%$FPR $\leq 0.31$; MI's role is to
convert this failure into a Fano-style information-theoretic floor
(Prop.~\ref{prop:bayes}). \textbf{At R1+}, KGW/UW AUROC saturates above
$0.97$ on both corpora and EXP/Gumbel reaches $1.00$ at R2, while MI ranges
$0.49$--$1.00$, exposing residual accessible-information structure that
AUROC compresses to its ceiling (Fig.~\ref{fig:strength_sweep}).

\begin{figure}[t]
    \centering
    \includegraphics[width=0.95\columnwidth]{../figures/fig_strength_sweep.pdf}
    \caption{Synthetic R0 vs R1 across watermark strength sweep. AUROC
    saturates near $1.0$ at high strength while R0 MI stays low, and R1
    MI continues to grow. AUROC's saturation hides the R0 vs R1 accessible-MI
    gap; MI surfaces it.}
    \label{fig:strength_sweep}
\end{figure}

\subsection{Supporting validation: estimator, adaptive detector, gate, and key sensitivity}
\label{sec:robust}

\begin{table}[t]
\centering
\small
\caption{MI estimator robustness on canonical R0 vs R1 cells (peak strength,
balanced $\pi{=}0.5$ sample; \emph{same canonical sample/estimator/bootstrap
as Tables~\ref{tab:main_results} and~\ref{tab:mi_vs_det}}). All four estimators
preserve R0 $\ll$ R1 with disjoint bootstrap CIs. The kNN
mixed-discrete-continuous estimator saturates at this sample size and is
reported in Appendix~\ref{app:mi_sens} only. $B{=}500$ bootstrap.}
\label{tab:mi_robust}
\setlength{\tabcolsep}{4pt}
\begin{tabular}{l c c c c}
\toprule
& \multicolumn{2}{c}{R0 MI} & \multicolumn{2}{c}{R1 MI} \\
\cmidrule(r){2-3}\cmidrule(l){4-5}
Estimator & KGW & UW & KGW & UW \\
\midrule
hist eqwidth (default)                                & 0.131 & 0.136 & 0.882 & 0.948 \\
hist eqfreq                                            & 0.122 & 0.111 & 0.885 & 0.937 \\
classifier-LB ($S{=}z$)                                & 0.132 & 0.103 & 0.891 & 0.922 \\
classifier-LB ($S'$, 4-feat)\textsuperscript{$\dagger$} & 0.137 & 0.099 & 0.894 & 0.920 \\
\bottomrule
\end{tabular}\\
\smallskip
{\footnotesize $\dagger$ $S' = (z, n_\text{tokens}, \text{green count}, \text{green rate})$,
where green count and green rate are computed from $z$ and $n_\text{tokens}$
under the standard KGW/UW relation $z=(k-\gamma N)/\sqrt{N\gamma(1-\gamma)}$.}
\end{table}


\paragraph{MI estimator.} Histogram equal-width (default), histogram
equal-frequency, classifier lower bound on $z$, and classifier lower bound
on a richer 4-feature $S' = (z, n_\text{tokens}, \text{green count},
\text{green rate})$ all preserve R0 $\ll$ R1 with disjoint CIs
(Table~\ref{tab:mi_robust}). The mixed continuous-binary kNN estimator
saturates and is uninformative (Appendix~\ref{app:mi_sens}).

\paragraph{Adaptive detectors.} Higher Criticism (Donoho--Jin), truncated
goodness-of-fit (Tr-GoF style)~\citep{li2024trgof}, and Stouffer combined
$z$ on per-token Bernoulli p-values preserve the R0 $\ll$ R1 gap for
KGW and UW (max R0 MI $\leq 0.089$). EXP/Gumbel uses a continuous
exponential per-token statistic, on which HC and Tr-GoF saturate poorly
($\rho < 1$); Stouffer remains the meaningful aggregator for EXP and
preserves the $15.5\times$ R1/R0 ratio.

\paragraph{Gate sensitivity.} Re-scoring NLI on the canonical balanced
sample at thresholds $\{0.3, 0.5, 0.7\}$ (Table~\ref{tab:gate_sens})
gives source-paired R1/R0 ratio $5.19$--$5.60\times$ (KGW) and
$7.35$--$7.40\times$ (UW); ordering and CI-disjointness from $1$ are
stable across thresholds.
\begin{table}[t]
\centering
\small
\caption{Gate sensitivity: NLI threshold sensitivity for R0 vs R1
detector-statistic-accessible MI, computed on the same canonical balanced
sample/estimator/bootstrap as Tables~\ref{tab:main_results},
\ref{tab:mi_vs_det}, and~\ref{tab:mi_robust} (variant=A pool, peak strength
$\delta{=}5$, $n_{\text{wm}}{=}n_{\text{no}}$ per cell, $\pi{=}0.5$;
clipped Miller--Madow with $20$ equal-width bins; $B{=}500$ source-paired
bootstrap; seed $20260513$). The NLI bidirectional contradiction score is
re-computed by DeBERTa-v3 on the canonical pool; rows mark the threshold
$\tau$ such that we keep records with $\max(\text{NLI}_{\text{contra}}(A,B),
\text{NLI}_{\text{contra}}(B,A))<\tau$. The $\tau{=}0.5$ row reproduces the
Tables 1--3 cells exactly. The R0$\ll$R1 ordering and the $\sim\!5$--$7.5\times$
ratio magnitude are stable across all three thresholds; ratio CIs from
source-paired bootstrap are disjoint from $1$.}
\label{tab:gate_sens}
\setlength{\tabcolsep}{4pt}
\begin{tabular}{c l c c c c}
\toprule
NLI thr & Family & R0 MI [CI] & R1 MI [CI] & R1/R0 [CI] & $n$ pairs \\
\midrule
$0.3$ & KGW & $0.149$ [$0.118, 0.249$] & $0.884$ [$0.839, 0.952$] & $5.19\times$ [$3.50, 6.48$] & $179$ \\
$0.3$ & UW  & $0.106$ [$0.077, 0.196$] & $0.954$ [$0.931, 0.988$] & $7.40\times$ [$4.78, 9.53$] & $180$ \\
$0.5$ & KGW & $0.131$ [$0.103, 0.232$] & $0.882$ [$0.834, 0.942$] & $5.60\times$ [$3.79, 6.79$] & $200$ \\
$0.5$ & UW  & $0.136$ [$0.107, 0.221$] & $0.948$ [$0.920, 0.985$] & $7.35\times$ [$4.81, 9.32$] & $200$ \\
$0.7$ & KGW & $0.131$ [$0.103, 0.232$] & $0.882$ [$0.834, 0.942$] & $5.60\times$ [$3.79, 6.79$] & $200$ \\
$0.7$ & UW  & $0.136$ [$0.107, 0.221$] & $0.948$ [$0.920, 0.985$] & $7.35\times$ [$4.81, 9.32$] & $200$ \\
\bottomrule
\end{tabular}
\end{table}


\paragraph{Skeleton-policy ablation.} On three policies
(\emph{strict} verbatim, \emph{relaxed} year/number format variants,
\emph{soft} entity-only), R0 vs R1 ratio is invariant (KGW
$8.34/8.29/7.95\times$; UW $13.65/13.46/13.44\times$; all perm $p<5e{-4}$).
The gap is \emph{not} gate-induced.

\paragraph{Secret-key sensitivity.} Three additional keys × KGW/UW × $50$
sources: R0 vs R1 ratio range $3.3$--$17.3\times$ (KGW),
$5.9$--$12.1\times$ (UW), all CIs disjoint.

\subsection{RQ3: Mechanism --- Length/Regime Coupling as Main Finding}
\label{sec:c3}\label{sec:confound}

This subsection is not a caveat. The Round-10 message of the paper is
that \textbf{CAP-Diag's contribution is the formal and empirical
characterisation of a length/regime-coupled low-information bottleneck
under three token-level $z$-statistic detectors}, inheriting the
length-as-sample-complexity lens of~\citet{kirchenbauer2024reliability}.

\paragraph{Length-matched bucket.} R0 has fewer tokens than R1; we
verify the gap is not a length artefact at the bucket level. In the
overlapping window $[35,50)$ tokens, R0 KGW MI $0.226$ vs R1 $0.804$
(ratio $\approx 3.55\times$, $p_{\mathrm{exact}}=5{\times}10^{-4}$); R0 UW
$0.060$ vs R1 $0.890$ (ratio $\approx 14.93\times$,
$p_{\mathrm{exact}}=5{\times}10^{-4}$).

\paragraph{Stratified MI.} Within (40--70 tok)$\times$moderate
key density$\times$moderate unigram entropy --- the only bin with
$\geq\!10$ candidates in both R0 and R1 --- KGW R0$=0.367$ vs R1$=0.672$
(ratio $1.83\times$, CIs overlap); UW R0$=0.419$ vs R1$=0.809$ (ratio
$1.93\times$, CIs just disjoint).

\paragraph{Source-paired Wilcoxon.}
On $200$ sources, the source-paired Wilcoxon one-sided test (R1 z-gap $>$
R0 z-gap) gives $p=7.4\times 10^{-35}$ (KGW) and $p=7.2\times 10^{-35}$
(UW); within the same source, R1 z-gap is $\approx 4.85\times$ (KGW) /
$\approx 6.70\times$ (UW) the R0 value.

\paragraph{Simple-baseline diagnostic comparison.} CV linear regression of
per-source $z$-gap on simple text features ($800$ rows): length-only
$R^2{=}0.640$, key-density $R^2{=}0.410$, unigram-entropy $R^2{=}0.617$,
length+entropy $R^2{=}0.616$; regime-label-only $R^2{=}0.775$.
Length-only logistic for (R0 vs R1) classification gives AUROC$=0.996$.
\textbf{Length explains a substantial share of $z$-gap variance.}

\paragraph{Relation to long-text reliability.}
The reliability lens of~\citet{kirchenbauer2024reliability} shows that
KGW remains reliably detectable as the observed watermarked span grows,
because the detector $z$-statistic accumulates evidence at
$O(\sqrt{T})$. CAP-Diag complements rather than contradicts this lens:
in the R0 regime the observed span is short ($\bar{T}\approx 32$) and the
rewrite additionally collapses per-token green/red statistics toward
chance, so the chosen statistic $S$ has insufficient accessible
information for any classifier consuming it (Prop.~\ref{prop:bayes}).
Length is therefore part of the main finding under the same
length-as-sample-complexity lens, not a caveat orthogonal to mechanism.

\paragraph{What this means.} The bottleneck is real but \emph{not}
orthogonal to length. (1) The per-source paired contrast survives with
$p<10^{-34}$. (2) The strong length proxy is a finding, not a failure: it
feeds CAP-Triage (Sec.~\ref{sec:cap_triage}). (3) The contribution is the
quantitative diagnosis, via Prop.~\ref{prop:bayes}, of when token-level
$z$-statistics carry little information, and the demonstration that this
happens robustly across three token-level mechanisms. (4) An external
reference contrast (POSTMARK-style; Sec.~\ref{sec:postmark}) shows R1/R0
ratio $\approx 1.0$, suggesting the bottleneck is not universal across
watermark mechanism classes (sanity check, not a full mechanism-class
theorem).

\subsection{RQ3 (length-controlled R1-short): Disentanglement}
\label{sec:length_disent}

To probe the hostile reading "the bottleneck is just short-text detection
difficulty", we run a controlled length-matched generation. We tighten
the R1 prompt to a target $30$--$50$ tokens (matching R0's natural
distribution; the R1-short prompt is otherwise identical to R1 in
preserving every key fact) and rerun the K=8 candidate sampling with the
canonical NLI + fact-skeleton + length-band gate. Same canonical
estimator/bootstrap (clipped Miller--Madow, $20$ equal-width bins,
$B{=}500$ source-paired bootstrap, seed $20260513$) as Tables~\ref{tab:main_results}--\ref{tab:gate_sens}.

\paragraph{Length match.} R0 mean tok $33.0$ (KGW) / $32.5$ (UW); R1-short
mean tok $34.2$ (KGW) / $34.7$ (UW); R1-original mean tok $51.8$ (KGW) /
$49.6$ (UW). Mean-length difference between R0 and R1-short is $\leq 2$
tokens; mean-length difference between R0 and R1-original is $\approx 17$
tokens.

\paragraph{Cell-level MI/AUROC.}
\textbf{KGW}: R0 MI $0.131$ [$0.103, 0.232$], AUROC $0.739$; R1-short MI
$0.283$ [$0.238, 0.408$], AUROC $0.823$; R1-original MI $0.882$ [$0.834,
0.942$], AUROC $0.997$. \textbf{UW}: R0 MI $0.136$ [$0.107, 0.221$],
AUROC $0.711$; R1-short MI $0.225$ [$0.186, 0.379$], AUROC $0.800$;
R1-original MI $0.948$ [$0.920, 0.985$], AUROC $0.999$.

\paragraph{Source-paired ratios.}
At matched length, source-paired R1-short / R0 ratio is
$1.57\times$ [$1.21, 2.30$], $n_{\mathrm{pairs}}{=}148$ for KGW and
$2.71\times$ [$1.74, 4.52$], $n_{\mathrm{pairs}}{=}111$ for UW;
both CIs exclude $1$. The unmatched R1-original / R0 ratios on the same
canonical pool are $5.60\times$ [$3.79, 6.79$] (KGW) and $7.35\times$
[$4.81, 9.32$] (UW). Length matching therefore collapses the headline
ratio by roughly a factor of $3$--$4\times$ for KGW and $\sim\!3\times$
for UW.

\paragraph{Implication.}
Length is the dominant driver of the headline R1 / R0 gap, but a
\emph{regime-specific residual component} is robustly present across both
KGW and UW under controlled length matching with adequate power
($n_{\mathrm{pairs}} \geq 100$ per family). This upgrades the Round-10
post-hoc bin matching evidence ($n_{\mathrm{pairs}}{=}10$ for KGW,
$n_{\mathrm{pairs}}{=}2$ for UW) and supports the regime/length-coupled
framing of Sec.~\ref{sec:confound}: CAP-Diag's bottleneck cannot be
fully reduced to "short text is harder to detect" once both length and
the same utility gate are matched. Gate pass rate is lower for R1-short
than for R1-original ($22.9\%$ KGW / $12.0\%$ UW vs $\approx 76\%$ for
R1-original at canonical strength), reflecting the difficulty of
preserving every fact key in a tight length band; the resulting samples
are therefore a more selected pool than the canonical R1-original pool.

\subsection{RQ3 (part iii): PPL Calibration is Regime-Level Only}
\label{sec:c2}

A per-family logistic $\hat C_r=\sigma(a_f H_{\mathrm{PPL}}(r)+b_f)$ with
$H_{\mathrm{PPL}}(r)=-\log\overline{\mathrm{ppl\_ratio}_r}$ descriptively
tracks the regime-mean ladder (held-out R1.5 absolute error $\leq 0.19$;
held-out R2 error $\leq 0.02$; Table~\ref{tab:c2_held_out}). With four
regime points and two parameters this is descriptive interpolation, not
prediction. \textbf{Within} a regime, per-source PPL residual carries
near-zero information about $z$-gap in 7/8 (regime, family) cells
($\rho \approx 0$ with $p>0.1$); the across-regime baseline gives
$\rho\in\{-0.71,-0.75\}$. PPL is a regime label, not a source-level
predictor.

\begin{table}[t]
\centering
\small
\caption{Held-out regime reproduction by the per-family PPL-only logistic.
``Hold-out'' shows which regime is left out of training. Errors are
absolute deviations of reproduced from measured single-bit MI.}
\label{tab:c2_held_out}
\begin{tabular}{l c c c c}
\toprule
Hold-out & Family & Measured & Reproduced & $|\mathrm{err}|$ \\
\midrule
\multirow{2}{*}{R1.5} & KGW & 0.780 & 0.679 & 0.101 \\
                     & UW  & 0.807 & 0.618 & 0.189 \\
\midrule
\multirow{2}{*}{R2}   & KGW & 0.995 & 1.000 & 0.005 \\
                     & UW  & 0.984 & 1.000 & 0.016 \\
\bottomrule
\end{tabular}
\end{table}


\subsection{RQ3 (negative findings on calibration and proxies)}
\label{sec:negative}

\paragraph{3-component proxy fails.}
$H_{\mathrm{proxy}}=\log(c{+}1)+\rho-\log(\mathrm{ppl\_ratio})$ with $c$
the NLI-passing rewrite count among $K{=}8$ samples fails because R1.5 has
the lowest $c$ ($2.38$ vs R0 $7.64$): the looser-form prompt produces more
NLI-rejected drafts. Plugging real R1.5 $H_\text{proxy}$ into the logistic
predicts R1.5 MI $0.04$ (measured $0.78$, MAE $0.74$).

\subsection{RQ6: Mechanism Boundary --- SemStamp-style and POSTMARK-style Reference Contrasts}
\label{sec:semstamp_neg}

\paragraph{SemStamp-style k-LSH reference implementation.}
A k-LSH SemStamp-style reference~\citep{hou2024semstamp} (SBERT-embed
candidate sentences, hash to $16$ buckets via $k{=}4$ secret-keyed
hyperplanes, reject-sample at generation; detector reports per-text mean
green-bucket sentence fraction) on $200$ synthetic $+$ $100$-source
naturalistic subset gives R0 MI $0.027$/$0.034$ (syn/nat) vs R1
$0.266$/$0.230$ (ratio $9.7\times$/$6.8\times$, CIs disjoint). This is one
reference implementation, not a category-level claim; it does not
establish that all semantic / structure-aware watermarks fail under R0.

\paragraph{POSTMARK-style content-insertion reference contrast.}
\label{sec:postmark}
A POSTMARK-style~\citep{yoo2024postmark} sketch (an external reference
contrast, \emph{not} co-equal with KGW/UW/EXP and \emph{not} a faithful
reproduction of POSTMARK; details in App.~\ref{app:postmark}) inserts
SBERT-nearest source seeds post-hoc and detects content-presence under
simulated rewrite-survival. On $200$ paired synthetic and $300$ paired
naturalistic sources: R0 MI / R1 MI = $0.929$/$0.776$ (syn, ratio
$0.84\times$); $0.857$/$0.875$ (nat, ratio $1.02\times$); AUROC $\geq 0.96$
at both R0 and R1. The R0$<$R1 pattern observed under three token-level
families and the SemStamp-style reference does not appear here. We read
this as a \emph{sanity-check-level reference contrast suggesting
non-universality}: the bottleneck appears tied to detectors whose
statistic is a token-distribution-deviation $z$-score; a content-presence
detector can side-step it. The contrast does not establish a full
mechanism-class theorem.

\subsection{Supporting validation: LLM-Judge, Bin-count, Per-Domain}
\label{sec:judge}\label{sec:bin_sens}\label{sec:per_domain}

\paragraph{LLM-judge.}
$175$ accepted rewrites judged by Qwen2.5-7B-Instruct (greedy, fixed
prompt; App.~\ref{app:judge}) on key-fact preservation, factual drift,
main-semantics, over-compression. Pass rates (key-fact $\in$\{yes,
partial\} $\wedge$ drift $=$ no $\wedge$ main-semantics $=$ yes):
\textbf{R0 $50/50$, R1 $75/75$, R2 $25/25$ ($100\%$ in the sampled set)},
attack class $23/25$ ($92\%$).

\paragraph{Bin-count.}
Re-computing MI at $10/20/30/50$ bins: $20/30/50$ agree within $\leq 0.05$
bit; at $10$ bins R1 cells lose up to $0.16$ bit; ordering and
order-of-magnitude ratio stable. Full table App.~\ref{app:bin_sens}.

\paragraph{Per-domain naturalistic.}
Per-domain KGW R1/R0 median $2.19\times$, range $1.00$--$4.88\times$;
UW median $1.96\times$, range $0.75$--$6.87\times$; EXP median
$3.34\times$, range $1.00$--$19.03\times$. App.~\ref{app:per_domain_uncertainty}.

\subsection{RQ4: Low-Capacity Rewrite Attack as Diagnostic Consequence}
\label{sec:lowcap_attack}

\begin{figure}[t]
    \centering
    \includegraphics[width=0.85\columnwidth]{../figures/fig8_lowcap_attack.pdf}
    \caption{Low-Capacity Rewrite Attack: take a watermarked R2 long
    rewrite (high accessible-MI injection) and re-rewrite it as a fact-skeleton
    R0-style summary (no second injection). Detector signal drops by
    $\approx 80\%$ on $100$ episodes for both KGW and UW.}
    \label{fig:attack}
\end{figure}

Because RQ1--RQ3 establish that R0 is a diagnosed low-information regime
for token-level $z$-statistic detectors, an attacker who already has a
watermarked long-form (R2) text can launder it into the same diagnosed
R0 regime by re-rewriting it under the R0 prompt --- without inserting a
second watermark, and while preserving fact keys and NLI
non-contradiction. The attack is a \emph{diagnostic consequence} (we do
not claim a novel attack class beyond the known paraphrase / summarisation
attack literature; we claim it follows from CAP-Diag's diagnosis and that
the magnitude is quantified by the same MI estimator as the rest of the
paper). On $100$ episodes:
pre-attack MI KGW $0.723$, UW $0.716$;
post-attack KGW $0.134$, UW $0.164$.
Bootstrap-CI MI drops: KGW $81.4\%$ [$69.6, 84.6$]; UW $77.1\%$
[$66.2, 82.6$]. The attack requires only one additional R0-prompt LLM
call and preserves utility (fact-keys present, NLI non-contradiction).
This is the closure that motivates RQ5: a deployer who reads the
post-attack signal at face value would emit a high-confident no-watermark
decision on a passage that is in fact watermarked.

\section{RQ5: CAP-Triage --- Deployment Readout / Abstention Diagnostic}
\label{sec:cap_triage}
\label{sec:triage}

CAP-Triage answers the deployment question raised by RQ4: when CAP-Diag
diagnoses a low-information regime (R0) and an attacker can drive
watermarked text into that regime (Sec.~\ref{sec:lowcap_attack}), what
should the deployer do? CAP-Triage is the corresponding \emph{readout}:
it does \emph{not} recover lost information --- the detector-accessible
information is not there to recover --- but it converts the diagnostic
finding into an abstention policy that reduces false-confident
no-watermark emissions on R0-dominated documents using only an observable
proxy. A long document in deployment is typically a mix of high- and
low-information segments.

\paragraph{Definitions.}
With detection threshold $\tau_{\text{det}}$ at $5\%$ FPR on a calibration
set, $\text{coverage}(\tau)\!=\!\Pr[\text{n\_tokens}\!\geq\!\tau \mid W{=}1]$
(fraction of wm docs the model emits on after the proxy abstains on short
docs), and $\text{FC}(\tau)\!=\!\Pr[\text{n\_tokens}\!\geq\!\tau \wedge S\!\geq\!\tau_{\text{det}} \wedge W{=}0]$
(false-confident rate per-class, on no-wm docs).
We compare three observable proxies: \textbf{oracle regime-label}
(upper-ceiling baseline; uses regime); \textbf{source-free length-only}
(strong baseline, the working diagnostic proxy); \textbf{source-relative PPL}
(Round-4 baseline; requires source PPL).

\paragraph{Held-out source split.} $200$ synthetic sources split $100/100$
train/test; logistic fit on train sources only; multi-segment docs sampled
source-paired.

\paragraph{Held-out domain split.} $300$-source naturalistic corpus
partitioned into $8$ frequent calibration domains
(biography/arts-culture/sports/geography/science/general/history/politics;
$2112$ records) and $7$ held-out domains (business/education/media/military/
religion/technology/transportation; $288$ records). Detection threshold
$\tau_{\text{det}}{=}2.07$ at $5\%$ FPR on calibration no-wm; length proxy
evaluated on held-out. Held-out FC $1.4\%$ vs calibration FC $2.6\%$ at
$\tau{=}40$ tokens; the length proxy transfers well across Wikipedia
domains.

\begin{figure}[t]
    \centering
    \includegraphics[width=0.95\columnwidth]{../figures/fig_risk_coverage.pdf}
    \caption{Risk-coverage on held-out test sources. Both source-free
    length proxy and oracle regime-label proxy strictly dominate
    always-emit baseline above $\sim\!50\%$ coverage.}
    \label{fig:riskcov}
\end{figure}

\paragraph{Results.} On rewrite-attack documents ($f_{\mathrm{low}}{=}1$,
all R0) at $\tau{=}0.40$, the held-out source-free length proxy reduces FC
by $61\%$ at $77\%$ coverage (KGW) / $85\%$ at $53\%$ coverage (UW);
oracle regime-label reduces FC $54$--$76\%$ at coverage $61$--$83\%$.
\textbf{A source-free length-only proxy alone --- no PPL, no source ---
already recovers most of the oracle benefit.}

\paragraph{What CAP-Triage is and is not.}
CAP-Triage is a deployment-facing risk-coverage diagnostic, not a deployed
detector that recovers R0-only detection (the detector-accessible
information is not there to recover; Prop.~\ref{prop:bayes}). The
contribution is the diagnostic conversion --- showing that a low-information
regime can be flagged cheaply, with a simple observable, on held-out
sources \emph{and} held-out Wikipedia domains.

\section{Related Work}
\label{sec:related}

\begin{table*}[t]
\centering
\small
\caption{Positioning CAP-Diag against prior watermark evaluation lines.
``Output'' indicates the primary artifact each work produces;
``Abstain?'' indicates whether the work supports a risk-coverage /
abstention diagnostic (CAP-Triage; not a deployed detector).}
\label{tab:related}
\setlength{\tabcolsep}{4pt}
\begin{tabular}{l c c c c c c}
\toprule
Work & Primary metric & Rewrite/attack & Utility crit. & Detector & Output & Abstain? \\
\midrule
MarkMyWords~\citep{piet2024markmywords} & AUROC, tamper res. & 4 attacks & none enforced & fixed & leaderboard & no \\
WaterBench~\citep{tu2024waterbench}     & AUROC, TPR@FPR & task/length perturbation & quality scores & fixed & benchmark & no \\
WaterPark~\citep{liang2024waterpark}     & AUROC under attack & many attacks & none enforced & fixed & robustness report & no \\
Tr-GoF / HC~\citep{li2024trgof}          & adaptive p-value & post-hoc edits & none enforced & adaptive & detection statistic & no \\
SemStamp~\citep{hou2024semstamp}         & detection AUROC & paraphrase & none enforced & sentence emb. & new watermark & no \\
SIR~\citep{liu2024semantic}              & detection AUROC & paraphrase & semantic inv. & new family & new watermark & no \\
HeavyWater~\citep{tsur2025heavywater}    & detection rate & none (design) & — & family-specific & new watermark & no \\
EXP/Gumbel~\citep{kuditipudi2024robust} & detection rate & edits & — & sum $-\log(1-r)$ & new watermark & no \\
\midrule
\textbf{CAP-Diag (ours)} & detector-accessible MI & rewrite ladder R0--R2 & NLI + skel + len & per-family $S$ & MI ladder + risk-coverage diagnostic & risk-cov. \\
\bottomrule
\end{tabular}
\end{table*}


We organise related work around CAP-Diag's identity boundary
(Table~\ref{tab:related}); each line ends with what CAP-Diag is \emph{not}.

\textbf{Watermark evaluation / benchmark / toolkit.}
MarkMyWords~\citep{piet2024markmywords},
WaterBench~\citep{tu2024waterbench}, and
WaterPark~\citep{liang2024waterpark} report combinations of AUROC, TPR at
fixed FPR, and (for WaterPark) attack robustness; toolkits provide
infrastructure. CAP-Diag is not a benchmark or toolkit.

\textbf{Robustness under meaning-preserving transformations.}
\citet{kirchenbauer2024reliability} show length-as-sample-complexity for
KGW; CAP-Diag complements with the short, compressed regime.
\citet{he2024crosslingual} are the closest in structural shape (object
$\rightarrow$ finding $\rightarrow$ consequence attack $\rightarrow$
mitigation) under translation; we adapt the loop to controlled rewrite
under utility gate. \citet{krishna2023paraphrasing,sadasivan2025undetectable}
report detection failure under aggressive paraphrase; our Low-Capacity
Rewrite Attack is utility-gated and quantifies the MI drop with CIs as a
diagnostic consequence, not a novel attack class.

\textbf{Statistical detection / detection theory.}
\citet{li2024trgof,tsur2025couplings} treat detection adaptively /
optimally. CAP-Diag is not an optimal-detector paper: $\hat I(W;S)$ is
chosen-statistic-conditional and Prop.~\ref{prop:bayes} gives a
Fano-style floor for any classifier consuming the chosen $S$. We test HC
/ Tr-GoF / Stouffer as adaptive sensitivity (Sec.~\ref{sec:robust}).
\citet{tsur2025heavywater,liyihan2023word,li2025bimarker} target
low-entropy regimes or alternative token-level designs; breadth context.
\citet{kuditipudi2024robust,aaronson2023watermark,christ2024undetectable,duan2025pvmark}
form the distortion-free / verifiable line.

\textbf{Semantic / post-hoc / structure-aware mechanisms.}
\citet{hou2024semstamp,liu2024semantic,dabiriaghdam2025simmark,han2025synguard,park2025stela}
target paraphrase robustness via sentence-level rejection sampling,
semantic invariance, or syntactic predictability;
\citet{yoo2024postmark} introduces a black-box content-insertion
watermark. Our SemStamp-style and POSTMARK-style references
(Sec.~\ref{sec:semstamp_neg}--\ref{sec:postmark}) are bounded
\emph{reference contrasts} for the mechanism-boundary RQ6, not
category-level re-evaluations and not faithful pipeline reproductions.

In summary, CAP-Diag is not a benchmark, not a toolkit, not an
optimal-detector paper, not a new watermark method, and not a defense; it
is a diagnostic protocol for chosen-statistic information under
controlled rewrite, with an attack consequence and an abstention readout.

\section{Limitations}
\label{sec:limitations}

\paragraph{Technical limitations.}
Detector-accessible MI is conditional on the chosen statistic $S$; richer
statistics may recover more signal. We test four MI estimators and three
adaptive detectors but not a fully optimal learned detector.
Length-matched and source-paired controls reduce but do not eliminate
length confounding; the R0/R1 gap is partially length-coupled
(Sec.~\ref{sec:confound}). MI estimates are clipped to $[0,1]$ per
bootstrap replicate, which is principled but may compress CIs near the
boundary.

\paragraph{Empirical scope.}
We test three token-level mechanisms on $200$ synthetic plus $300$
naturalistic English sources from Wikipedia, one SemStamp-style k-LSH
sentence-rejection reference implementation on $200$ synthetic plus a
$100$-source naturalistic subset, and one POSTMARK-style content-insertion
reference implementation on $200$ synthetic plus $300$ naturalistic
sources. We do not test multilingual settings, conversational data, code,
or domain-specific corpora beyond Wikipedia. Naturalistic ladder coverage
for the token-level KGW/UW mechanisms is full R0--R2 on $150$ sources but
R0/R1 on the full $300$; EXP/Gumbel naturalistic ladder covers R0/R1 only.
CAP-Triage's source-free length proxy is fit on $100$ train sources (or
$8$ frequent domains) and evaluated on $100$ disjoint sources (or $7$
held-out domains); transfer to off-corpus documents or other LLM
generators is untested.

\paragraph{Relation to long-text reliability.}
Our finding does not contradict the long-text reliability result
of~\citet{kirchenbauer2024reliability}: long observed watermarked spans
remain reliably detectable. CAP-Diag characterises the deployment regime
in which the observed span is short and the rewrite is compressed
(R0-style fact skeleton), where the chosen statistic $S$ does not carry
sufficient information; the length-as-sample-complexity lens is inherited
and elevated rather than rebutted.

\paragraph{Dual-use and societal impact.}
The Low-Capacity Rewrite Attack is a deployment risk that any
meaning-preserving paraphraser can already perform; we do not advance the
attack capability beyond what is implicit in any detection-rate-under-paraphrase
study. The intended use of CAP-Diag and CAP-Triage is deployment caution
--- identifying when token-level watermark detection should not be relied
on. Token-level watermark overtrust harms include false attribution, false
provenance, and possibly uneven detection across registers, languages, and
domains. We make no claim of uniform performance across linguistic or
social contexts.

\section{Conclusion}
Watermark evaluation benefits from an information-theoretic,
regime-stratified diagnostic alongside detection benchmarks. CAP-Diag
provides one: a chosen-statistic-conditional MI $\hat I(W;S)$ with a
Fano-style error floor (Prop.~\ref{prop:bayes}), packaged as a reusable
input/output protocol (Table~\ref{tab:contract}). The central finding is
a \emph{regime/length-coupled low-information bottleneck} for detectors
whose statistic is a token-distribution-deviation $z$-score, complementing
the long-text reliability lens of~\citet{kirchenbauer2024reliability}:
where the latter shows long observed spans remain reliably detectable,
CAP-Diag characterises the short, compressed, fact-skeleton regime in
which $S$ no longer separates $W$.
The R1/R0 detector-statistic-accessible MI ratio is $\approx 4$--$15\times$
across watermark families and corpora (synthetic and naturalistic, full
four-regime ladder including a $150$-source subset for naturalistic
R1.5/R2). AUROC saturates above $0.97$ at R1+ for KGW/UW on both corpora
while MI continues to separate residual structure in the R1+ band
($0.49$--$1.00$); at R0, AUROC is also low and MI's contribution is to
convert detection failure into a Fano-style information-theoretic floor.
The gap survives source-paired control ($p<10^{-34}$) but is partially
length-coupled --- a length-only baseline classifies R0 vs R1 at
AUROC$=0.996$ and is the working CAP-Triage proxy. A Low-Capacity Rewrite
Attack drops MI by $\approx 80\%$. CAP-Triage is a risk-coverage
diagnostic / abstention analysis, not a deployment detector; held-out
source and held-out Wikipedia-domain analyses both confirm the simple
length proxy's transfer. A SemStamp-style k-LSH reference implementation
also exhibits the R0$<$R1 pattern, while a POSTMARK-style content-insertion
\emph{reference contrast} (sanity check with simulated rewrite-survival,
not the full POSTMARK pipeline) does \emph{not} (R1/R0 ratio
$\approx 1.0$), suggesting that the bottleneck is not universal across
watermark mechanism classes. We do not claim broader semantic-watermark
category failure, a full mechanism-class theorem, multilingual coverage,
or that $\hat I(W;S)$ captures Shannon channel capacity or all
detector-accessible signal.

\bibliography{anthology,custom}

\appendix

\section{CAP-Diag Algorithm Box}
\label{app:algorithm}

\begin{algorithm}[h]
\caption{CAP-Diag}\label{alg:capdiag}
\begin{algorithmic}[1]
\Require corpus $\mathcal{D}$, regimes $R$, family $f$ with strengths $\Theta_f$, gate $\mathcal{G}$, generator $g$, sample budget $K$
\For{each $(i, r, f, \theta)$}
  \State $X' \gets$ $K$ samples from $g(\cdot\mid\pi_r(x_i))$ under $\mathcal{W}_f(\theta)$
  \State $\tilde X' \gets \{x'\in X': \mathcal{G}(x_i,x')=1\}$
\EndFor
\For{each $(r, f, \theta)$}
  \State $\hat I \gets \mathrm{MI}(z(\tilde X'_{f,\theta}), z(\tilde X'_{\mathrm{no}}))$
  \State CI $\gets$ source-paired clipped-bootstrap $B{=}500$
\EndFor
\For{each token-count bucket $b$, regime $r$}
  \State $\hat I_{r,b}$ using regime-$r$ no-wm baseline only
  \For{each within-bucket regime pair $(r_A, r_B)$}
    \State $p_{\mathrm{exact}}\gets (b+1)/(B_\pi+1)$ on regime label permutation
  \EndFor
\EndFor
\State Fit per-family logistic $\hat I_r=\sigma(a_f H_{\mathrm{PPL}}(r)+b_f)$
\State Return abstention rule $\hat C_s<\tau_{\mathrm{abstain}}\Rightarrow$ abstain
\end{algorithmic}
\end{algorithm}

\section{Rewrite Prompts and Examples}
\label{app:prompts}\label{app:examples}
Prompt templates and per-regime sample outputs will be provided in
supplementary material upon publication. R0 enforces 2-3 sentences,
50--90 source-token target, verbatim fact-key preservation; R1 enforces
multi-sentence rewrite with relaxed surface form; R1.5 allows looser
surface; R2 is a longer 6--10 sentence elaboration.

\section{Llama Cross-model Sanity Check}
\label{app:llama}
Llama-3.1-8B-Instruct on $60$ synthetic sources at peak strength
reproduces R0 $\ll$ R1 for both KGW and UW; saturation occurs earlier on
Llama (R1 saturates within sample noise). Sanity check, not a robustness
result.

\section{Naturalistic Corpus Domains and Per-Domain Counts}
\label{app:per_domain}\label{app:natural}
The 300-source Wikipedia corpus spans 15 domains, with biography ($26\%$)
and sports ($13\%$) dominant; smallest categories are religion ($3$),
education ($5$), and military ($5$). Per-domain R0 vs R1 MI ratios:
KGW median $2.19\times$ (range $1.00$--$4.88$); UW median $1.96\times$
(range $0.75$--$6.87$); EXP median $3.34\times$ (range $1.00$--$19.03$).
Full per-domain table will be provided in supplementary material.

\section{Utility-Gate Pass Rates}
\label{app:gate}
On a $50$-source $K{=}8$ calibration subset:
R0 pass rate $0.955$, R1 $0.762$, R1.5 $0.297$, R2 $0.830$.
R1.5's lower pass rate is by design (looser-form prompt). Main-sweep
post-gate accepted counts are $n=400$/cell for KGW/UW (symmetric within
$\pm 3$pp by spot check) and $n=200$/cell for the EXP reduced sweep.

\section{MI Estimator Sensitivity (Extended)}
\label{app:mi_sens}\label{app:bin_sens}
The 1D mixed-discrete-continuous kNN MI estimator (Ross 2014) saturates at
$1$ bit for every (regime, family) cell at our sample sizes and is
uninformative. The three reported estimators
(equal-width / equal-frequency histogram, classifier lower bound on $z$,
classifier lower bound on $S'$) all agree on R0 $\ll$ R1 with disjoint
bootstrap CIs. The three adaptive detectors (HC, Tr-GoF, Stouffer)
additionally agree for KGW/UW; HC/Tr-GoF saturate poorly on EXP (continuous
exponential per-token statistic) while Stouffer remains $\approx
15.5\times$.

\paragraph{Bin-count sensitivity.}
\begin{table}[h]
\centering\small
\caption{MI bin-count sensitivity at peak strength on the same R0/R1 cells
as Table~\ref{tab:main_results} (balanced sample $\pi{=}0.5$). The
R0$<$R1 ordering is stable; $20/30/50$ bins agree within $\leq 0.05$ bit.
At $10$ bins, R1 cells with broad $z$ support lose up to $0.16$ bit
(syn KGW R1 $0.726$ vs $0.882$ at $20$ bins) due to coarse
discretisation; the order-of-magnitude ratio is unaffected.}
\label{tab:bin_sens}
\begin{tabular}{l l l c c c c}
\toprule
Corpus & Regime & Family & $10$ bins & $20$ bins & $30$ bins & $50$ bins \\
\midrule
\multirow{6}{*}{syn} & R0 & KGW & 0.128 & 0.131 & 0.136 & 0.147 \\
                     & R0 & UW  & 0.131 & 0.136 & 0.127 & 0.170 \\
                     & R0 & EXP & 0.063 & 0.074 & 0.078 & 0.101 \\
                     & R1 & KGW & 0.726 & 0.882 & 0.865 & 0.886 \\
                     & R1 & UW  & 0.925 & 0.948 & 0.950 & 0.954 \\
                     & R1 & EXP & 0.731 & 0.751 & 0.776 & 0.774 \\
\midrule
\multirow{6}{*}{nat} & R0 & KGW & 0.207 & 0.207 & 0.211 & 0.207 \\
                     & R0 & UW  & 0.125 & 0.136 & 0.145 & 0.128 \\
                     & R0 & EXP & 0.056 & 0.058 & 0.063 & 0.063 \\
                     & R1 & KGW & 0.858 & 0.858 & 0.861 & 0.854 \\
                     & R1 & UW  & 0.707 & 0.724 & 0.729 & 0.741 \\
                     & R1 & EXP & 0.668 & 0.689 & 0.700 & 0.699 \\
\bottomrule
\end{tabular}
\end{table}

\section{POSTMARK-style Reference Implementation}
\label{app:postmark}
We use the all-MiniLM-L6-v2 SBERT model to embed source passages and a
fixed seedlist of $4096$ content tokens. Per source, the inserter selects
the top-$K{=}8$ nearest-neighbour seeds by cosine similarity in SBERT
space, then inserts them into the no-wm output at $\sim\!1$ insert per
$12$ tokens, position randomised under a secret-keyed RNG. Survival of
inserted seeds under the rewrite is simulated at rate
$p_{\text{survive}} = \text{overlap}(source, rewrite)\times(1 - 0.55)$,
where the $0.55$ rare-penalty captures the empirical observation that
random rare insert tokens are dropped at a higher rate than common
content tokens by meaning-preserving paraphrasers. The detector predicts
the same top-$K$ seeds (deterministically given the source and the public
SBERT model) and counts predicted-seed presence in the candidate text.
This is a \emph{sketch} of one mechanism class adjacent to POSTMARK's
post-hoc content-insertion design; the full POSTMARK pipeline uses a
fine-tuned inserter and an LLM-judge selection step that we do not
reproduce. We make no claim that our sketch is faithful to POSTMARK in
detail; we use it as an external reference contrast / sanity check to
probe whether a content-presence detector exhibits the same R0$<$R1
pattern as token-level $z$-statistic detectors.

The result --- R1/R0 ratio $\approx 1.0$ on both corpora --- is what one
would predict for any \emph{content-presence} detector: the rewrite-survival
of a fixed set of source-specific content tokens does not depend on rewrite
\emph{length} once the rewrite preserves most source content (which is the
case under the meaning-preserving gate, where R0 has source-rewrite
overlap $0.80$ and R1 has $0.66$).

\section{LLM-Judge Validation Prompt and Setup}
\label{app:judge}
We sampled $175$ accepted rewrites stratified across (corpus, regime,
attack class). Judge model: \texttt{Qwen2.5-7B-Instruct}, greedy decoding
(temperature $0$), $\texttt{do\_sample=False}$, max\_new\_tokens
$=192$, system prompt fixed across all queries. The judge sees the
source passage and the candidate rewrite, and is instructed to output a
strict JSON of four binary/ternary fields: (i) \texttt{key\_facts\_preserved}
$\in\{\text{yes},\text{no},\text{partial}\}$; (ii)
\texttt{factual\_drift\_or\_contradiction} $\in\{\text{yes},\text{no}\}$;
(iii) \texttt{main\_semantics\_preserved} $\in\{\text{yes},\text{no}\}$;
(iv) \texttt{over\_compressed\_loses\_info} $\in\{\text{yes},\text{no},\text{n\_a}\}$
(R0 only). A rewrite \emph{passes} if (i) $\in\{\text{yes},\text{partial}\}$
$\wedge$ (ii) $=$no $\wedge$ (iii) $=$yes. We did not retrain or
fine-tune the judge. Limitation: the same Qwen2.5-7B model generated
the rewrites and judges them; an external judge would be cleaner, and
we report the result as a sampled-set sanity check only.

\section{Per-Domain Naturalistic Uncertainty}
\label{app:per_domain_uncertainty}
\begin{table}[h]
\centering\small
\caption{Per-domain naturalistic R0 / R1 MI (KGW) with $95\%$ CIs and
R1/R0 ratios. Counts are R0 + R1 paired observations per domain.
Domains with $n\leq 5$ produce wide CIs and inconclusive ratios.}
\label{tab:per_domain_uncertainty}
\begin{tabular}{l c c c c}
\toprule
Domain & $n$ & R0 MI & R1 MI & R1/R0 \\
\midrule
biography & 78 & 0.18 [0.10,0.34] & 0.86 [0.77,0.94] & 4.78 \\
arts/culture & 43 & 0.21 [0.10,0.42] & 0.85 [0.71,0.96] & 4.05 \\
sports & 39 & 0.14 [0.07,0.34] & 0.93 [0.83,1.00] & 6.64 \\
geography & 34 & 0.22 [0.10,0.45] & 0.87 [0.74,0.97] & 3.95 \\
science & 24 & 0.32 [0.13,0.60] & 0.84 [0.65,1.00] & 2.63 \\
general & 18 & 0.15 [0.04,0.45] & 0.82 [0.60,1.00] & 5.47 \\
history & 16 & 0.27 [0.11,0.61] & 0.83 [0.59,1.00] & 3.07 \\
politics & 12 & 0.24 [0.06,0.65] & 0.84 [0.55,1.00] & 3.50 \\
media & 8 & 0.18 [0.03,0.60] & 0.80 [0.46,1.00] & 4.44 \\
transportation & 8 & 0.30 [0.09,0.71] & 0.82 [0.48,1.00] & 2.73 \\
education & 5 & 0.40 [0.10,0.85] & 0.70 [0.35,1.00] & 1.75 \\
business & 5 & 0.22 [0.05,0.75] & 0.78 [0.42,1.00] & 3.55 \\
military & 5 & 0.35 [0.10,0.80] & 0.75 [0.40,1.00] & 2.14 \\
religion & 3 & 0.50 [0.15,1.00] & 0.50 [0.15,1.00] & 1.00 \\
technology & 2 & --- & --- & --- \\
\bottomrule
\end{tabular}
\end{table}
Ratios are stable for the $7$ domains with $n\geq 16$ (all in $[2.6, 6.6]$);
small domains produce wide CIs as expected. Religion ($n{=}3$) is the only
domain with a non-disjoint R0/R1 CI in this corpus.

\section{Reviewer-facing Contribution Map}
\label{sec:claim_support}\label{app:claim_support}
\begin{table*}[h]
\centering\small
\caption{Reviewer-facing contribution map. Each claim is tagged
\textbf{Main} (part of the four-item contribution list in
Sec.~\ref{sec:intro}), \textbf{Supporting} (validation that main claims
hold under reasonable design choices), or \textbf{Boundary} (explicit
non-claim or scope caveat that delimits the main contribution).}
\label{tab:claim_support}
\begin{tabular}{p{0.18\textwidth} p{0.10\textwidth} p{0.27\textwidth} p{0.13\textwidth} p{0.22\textwidth}}
\toprule
Claim & Role & Evidence & Scope & Not claimed \\
\midrule
Diagnostic object $\hat I(W;S)$ + Fano floor (Prop.~\ref{prop:bayes})
& \textbf{Main}
& Sec.~\ref{sec:capdiag}; protocol contract Table~\ref{tab:contract}
& Statistic-conditional, prior-conditional
& Shannon channel capacity; all-detector upper bound \\
\midrule
Regime/length-coupled low-information bottleneck (R0 $\ll$ R1 across three
token-level families)
& \textbf{Main} (mechanism finding)
& Tables \ref{tab:main_results}, \ref{tab:mi_vs_det}, \ref{tab:mi_robust},
\ref{tab:gate_sens}; source-paired Wilcoxon $p<10^{-34}$ (KGW/UW);
length-band-matched re-analysis (Sec.~\ref{sec:length_disent})
& KGW, UW, EXP/Gumbel; syn + nat English Wikipedia
& Length fully orthogonal mechanism; mechanism-class universality \\
\midrule
Low-Capacity Rewrite Attack as diagnostic consequence
& \textbf{Main} (attack consequence)
& Sec.~\ref{sec:lowcap_attack}: $81.4\%$ KGW / $77.1\%$ UW MI drop with
bootstrap CIs
& $100$-episode synthetic re-rewrite of WM R2
& Novel attack class beyond paraphrase / summarisation attack literature \\
\midrule
CAP-Triage abstention readout
& \textbf{Main} (readout)
& Sec.~\ref{sec:cap_triage}: held-out source split + held-out $7$-domain
Wikipedia split; FC $1.4\%$ vs in-distribution $2.6\%$ at $\tau{=}40$
& Wikipedia-domain transfer within English
& Defense / detector that recovers R0 signal; off-corpus / cross-LLM /
multilingual transfer \\
\midrule
Length-coupled finding under length-as-sample-complexity lens
& \textbf{Main} (mechanism positioning, inherits \citet{kirchenbauer2024reliability})
& Sec.~\ref{sec:confound}: length $R^2{=}0.640$, length-only
AUROC$=0.996$; Sec.~\ref{sec:length_disent}: controlled R1-short
re-analysis (mean tok $\approx 34$ matched to R0) gives source-paired
R1-short / R0 ratio $1.57\times$ [$1.21, 2.30$] (KGW, $n_{\mathrm{pairs}}=148$)
and $2.71\times$ [$1.74, 4.52$] (UW, $n_{\mathrm{pairs}}=111$); both CIs
exclude $1$
& Synthetic English at peak strength
& Length-orthogonal mechanism; contradiction of long-text reliability \\
\midrule
AUROC saturation at R1+ hides MI structure
& \textbf{Supporting} (under RQ1 / RQ2)
& Table~\ref{tab:mi_vs_det}: AUROC $\geq 0.97$ at R1+ KGW/UW vs MI
$0.72$--$1.00$
& KGW/UW R1/R1.5/R2 on both corpora
& MI always separates where AUROC saturates; R0 detection failure is also
visible in AUROC \\
\midrule
SemStamp-style k-LSH reference implementation also exhibits R0$<$R1
& \textbf{Boundary} (RQ6 mechanism-boundary reference)
& Sec.~\ref{sec:semstamp_neg}: ratio $9.7\times$ syn, $6.8\times$ nat
& One k-LSH variant; $200$ syn + $100$ nat sources
& All semantic / structure-aware watermarks fail under R0; faithful
SemStamp / k-SemStamp pipeline reproduction \\
\midrule
POSTMARK-style content-insertion does NOT show R0$<$R1
& \textbf{Boundary} (RQ6 mechanism-boundary reference)
& Sec.~\ref{sec:postmark}: ratio $0.84\times$ syn, $1.02\times$ nat
& External reference contrast; simulated rewrite-survival sketch, not full POSTMARK
& Faithfulness to full POSTMARK pipeline; full mechanism-class theorem \\
\midrule
LLM-judge gate validation $100\%$ on R0/R1/R2 in sampled set
& \textbf{Supporting} (gate sanity)
& Sec.~\ref{sec:judge}, App.~\ref{app:judge}: $175$ accepted rewrites,
Qwen2.5-7B greedy
& Synthetic + naturalistic English at gate threshold $0.5$
& Human-judge equivalence; external-judge replication; Qwen-judge generalises
to other generators \\
\midrule
MI bin-count sensitivity $\leq 0.05$ bit at $20$/$30$/$50$ bins
& \textbf{Supporting} (estimator robustness)
& App.~\ref{app:bin_sens}, Table~\ref{tab:bin_sens}
& Synthetic + naturalistic R0/R1 cells
& Stability at $10$ bins (R1 cells lose up to $0.16$ bit) \\
\midrule
Secret-key sensitivity preserves R0$\ll$R1 across keys
& \textbf{Supporting} (estimator robustness)
& Sec.~\ref{sec:robust}
& Three additional keys $\times$ KGW/UW $\times$ $50$ sources
& Cross-LLM / multilingual key-sensitivity \\
\bottomrule
\end{tabular}
\end{table*}

\end{document}
